% !TEX program = xelatex
\documentclass[profile=ieee,lang=en]{skedreport}

\skedtitle{Software Requirements Specification for a Personal Media Content Tracking Information System}
\skedauthor{Victoria Buksha}
\skedaffiliation{Student, Information Systems Design Course}
\skedorganization{}
\skeddoctype{Software Requirements Specification}
\skedversion{1.3.0}
\skeddate{September 25, 2026}
\skedudc{004.415.2}
\skedkeywords{software requirements specification, media tracker, web application, SRS, content catalog, Socratic dialogue}
\skedplatform{Web browser}

\begin{document}

\skedmaketitle

\begin{abstract}
This document is a Software Requirements Specification (SRS) for the personal media content tracking web application covering books, TV series, and films. The system is intended for registered users who maintain a personal catalog of viewed and read content. The document defines functional requirements (31~requirements), non-functional requirements (9~requirements), acceptance criteria, BDD scenarios, and a traceability matrix for version~v1 (MVP). The document was prepared using the SKED methodology (Socratic Knowledge Elicitation \& Documentation) as part of the ``Information Systems Design'' course. The MVP includes: registration and authentication, a catalog with CRUD operations, tags with autocomplete, a Sessions mechanism for recording repeated viewings, collections (many-to-many), automatic completion date fill, and personal statistics. Social interaction features, offline mode, and integration with third-party media databases are deferred to future versions.
\end{abstract}

\skedtoc

\section{Introduction}

\subsection{Purpose and Scope of the Document}

This document is a Software Requirements Specification (SRS) for the personal media content tracking information system. The purpose of this document is to formally describe the functional and non-functional requirements of the system, define the boundaries of version~v1 (MVP), and provide a basis for software design documentation (SDD) and testing. \sked{K10}

\subsection{System Purpose}

The system is a web application that allows a registered user to maintain a personal catalog of media content --- books, TV series, and films. The application provides tracking of viewed and read content: recording statuses, ratings, notes, tags, repeated viewings (Sessions), organizing items into collections, and viewing personal statistics. \sked{K2,K4,K3,K5}

\subsection{Document Context and Audience}

This document was prepared as part of the ``Information Systems Design'' course (SDD cycle). Target audience: the course instructor and system developers. A reader of this document should obtain sufficient information to develop the system architecture and testing plan without additional clarification. \sked{K44}

\subsection{Document Structure}

\begin{itemize}
  \item \secref{sec:zahalnyi-opys} --- general system description, MVP boundaries, user roles, operating environment.
  \item \secref{sec:hlosarii} --- key term definitions.
  \item \secref{sec:funktsionalni-vymohy} --- complete list of functional requirements grouped by functional blocks.
  \item \secref{sec:nefunktsionalni-vymohy} --- non-functional requirements: performance, security, reliability, usability, compatibility, validation.
  \item \secref{sec:kryterii-pryiniattia} --- acceptance criteria for each requirement.
  \item \secref{sec:bdd} --- BDD scenarios in Gherkin format.
  \item \secref{sec:matrytsia} --- requirements traceability matrix.
  \item Appendix~A --- features outside the MVP scope.
\end{itemize}

\section{General System Description}
\label{sec:zahalnyi-opys}

\subsection{Product Overview and Perspective}

The system is a standalone web application with no external service integrations within the MVP scope. Each registered user has their own isolated data space --- a catalog, collections, and statistics. The system is not part of a larger platform and does not provide a public API in version~v1. \sked{K38,K41,K42}

\subsection{Functional Blocks}

The system implements the following functional blocks:
\begin{enumerate}
  \item Authentication and account management.
  \item Personal media content catalog.
  \item Content item management (CRUD).
  \item Tags and tag-based filtering.
  \item Sessions --- recording repeated viewings or readings.
  \item Collections --- custom named groups of items.
  \item Automatic completion date fill.
  \item Personal statistics.
\end{enumerate}
\sked{K10,K13,K16,K18,K21,K27,K28}

\subsection{User Roles}

Version~v1 has a single user role. No administrative role exists.

\begin{table}[h]
\caption{System User Roles}
\label{tab:roli}
\begin{tabular}{L{3cm} L{10cm}}
\toprule
\textbf{Role} & \textbf{Description} \\
\midrule
User & Registered account owner. Has access exclusively to their own catalog, collections, and statistics. May perform all CRUD operations on their own data. \\
\bottomrule
\end{tabular}
\end{table}
\sked{K44}

\subsection{Operating Environment}

The system is implemented as a web application. A native mobile application is not planned for version~v1. The system must function correctly in the latest two versions of Chrome, Firefox, Edge, and Safari browsers, including mobile Safari on iOS. The backend implementation language is not specified by this document and is determined at the architecture design stage. \sked{K38,K39,K45}

\subsection{Design and Implementation Constraints}

\begin{itemize}
  \item Implementation exclusively as a web application.
  \item Version~v1 supports only one interface language.
  \item The system does not integrate with external media content databases (TMDB, Google Books, etc.).
  \item Each user has access exclusively to their own data; cross-user access is not provided.
\end{itemize}
\sked{K38,K40,K41,K42}

\subsection{MVP Boundaries --- What Is Not Included in Version~v1}

The following features are technically feasible but deliberately excluded from the MVP to constrain the scope of the first release. A detailed list is provided in Appendix~A. \sked{K43,K50}

\section{Glossary}
\label{sec:hlosarii}

\Tabref{tab:hlosarii} lists the definitions of key terms used in this document.

\begin{table}[h]
\caption{System Glossary}
\label{tab:hlosarii}
\begin{tabular}{L{3.5cm} L{3.5cm} L{7cm}}
\toprule
\textbf{Term} & \textbf{Ukrainian equivalent} & \textbf{Definition} \\
\midrule
Catalog & Каталог & The complete list of all content items added by the user. \sked{K1} \\
\addlinespace
Content Item & Елемент контенту & A single entry in the catalog --- a book, TV series, or film with its own attributes (type, status, rating, tags, etc.). \sked{K2} \\
\addlinespace
Collection & Добірка & A user-created named group of arbitrary catalog items; one item can belong to multiple collections simultaneously. \sked{K3} \\
\addlinespace
Session & Session & An individual record of a viewing or reading of an item; contains date, rating, and note; one item can have multiple Sessions. \sked{K4} \\
\addlinespace
Tag & Тег & An arbitrary text label that the user attaches to an item for categorization. \sked{K5} \\
\addlinespace
Status & Статус & The current state of work on an item: \emph{planned~/ in progress~/ completed~/ on hold}. \sked{K6} \\
\bottomrule
\end{tabular}
\end{table}

\section{Development Methodology}
\label{sec:metodolohiia}

\subsection{Specification-Driven Development (SDD)}

The system is developed in accordance with the Specification Driven Development (SDD) methodology. The specification is the controlling source of system behavior --- no source code may define behavior independently of this document. \sked{K53}

Required process for every behavioral change:

\begin{enumerate}
  \item Update this document before making any code changes.
  \item Assign or update a requirement identifier (REQ-F-NNN or REQ-NF-NNN).
  \item Define an acceptance criterion (AC-NNN) for each requirement.
  \item Write an automated test before implementation (TDD, \secref{sec:metodolohiia}).
  \item Run the test and confirm it fails for the expected reason (Red).
  \item Implement the minimal code that passes the test (Green).
  \item Refactor only after tests pass (Refactor).
  \item Update log artifacts after completing the iteration.
\end{enumerate}

\subsection{Mandatory Test-Driven Development (TDD)}

TDD is mandatory for all implementation work. The Red~/ Green~/ Refactor cycle applies to every functional change: \sked{K53}

\begin{itemize}
  \item \textbf{Red} --- write a failing test first.
  \item \textbf{Green} --- implement the minimal code that passes the test.
  \item \textbf{Refactor} --- improve the code without changing external behavior.
\end{itemize}

No production code is accepted unless it is covered by at least one automated test. Real external systems (DBMS, third-party APIs) are not used in unit tests --- deterministic substitutes (stubs~/ mocks) are used instead.

\subsection{Traceability Rule}

Every implemented system behavior must have full traceability in the chain:

\begin{codelisting}
\caption{Traceability Chain}
\label{lst:trace}
\begin{skedcode}
REQ -> AC -> Test -> Implementation
\end{skedcode}
\end{codelisting}

Functionality without a requirement identifier is not part of the system. A requirement without an acceptance criterion is not ready for implementation.

\subsection{Test Log Artifacts}

After each test run, the development agent records a log file with the results. \sked{K54} The log file must contain:

\begin{itemize}
  \item total number of tests and the count of passed~/ failed;
  \item the status of each individual test (passed~/ failed~/ skipped);
  \item execution time of the test suite.
\end{itemize}

The log file is a machine-readable verification artifact --- based on it, the development agent and human reviewer determine whether the task has been completed according to the specification (the ``Is~verified?'' step in the SDLC process).

\section{Functional Requirements}
\label{sec:funktsionalni-vymohy}

\subsection{Authentication and Account Management}

\begin{table}[h]
\caption{Authentication and Account Requirements}
\label{tab:auth}
\begin{tabular}{L{2.5cm} L{11.5cm}}
\toprule
\textbf{ID} & \textbf{Requirement} \\
\midrule
REQ-F-001 & The system shall allow user registration with email, password, and nickname (all three fields are mandatory). Re-registration with the same email is not allowed. \sked{K7} \\
\addlinespace
REQ-F-002 & The system shall allow account login with email and password. \sked{K8} \\
\addlinespace
REQ-F-003 & The system shall allow account logout. After logout, access to protected pages requires re-authentication. \sked{K9} \\
\addlinespace
REQ-F-029 & The user may change their nickname in account settings. \sked{K49} \\
\bottomrule
\end{tabular}
\end{table}

\subsection{Catalog}

\begin{table}[h]
\caption{Catalog Requirements}
\label{tab:kataloh}
\begin{tabular}{L{2.5cm} L{11.5cm}}
\toprule
\textbf{ID} & \textbf{Requirement} \\
\midrule
REQ-F-004 & The system shall display a personal catalog --- a list of all user content items. \sked{K10} \\
\addlinespace
REQ-F-005 & The system shall provide catalog filtering by item type (book~/ TV series~/ film). \sked{K11} \\
\addlinespace
REQ-F-006 & The system shall provide catalog filtering by item status. Available status values: ``planned'', ``in progress'', ``completed'', ``on hold''. \sked{K12} \\
\addlinespace
REQ-F-023 & The system shall provide catalog item search by title. \sked{K29} \\
\addlinespace
REQ-F-026 & The system shall provide catalog sorting: by date added (newest~/ oldest) and by rating (highest to lowest and vice versa). \sked{K32} \\
\bottomrule
\end{tabular}
\end{table}

\subsection{Content Item Management}

\begin{table}[h]
\caption{Content Item Management Requirements}
\label{tab:content}
\begin{tabular}{L{2.5cm} L{11.5cm}}
\toprule
\textbf{ID} & \textbf{Requirement} \\
\midrule
REQ-F-007 & The system shall allow adding a new content item with attributes: title (required), type (required), status, rating~(0--10), notes, tags. The form displays type-specific optional fields (REQ-F-024) and a cover image field (REQ-F-025). \sked{K13} \\
\addlinespace
REQ-F-008 & The system shall allow editing any attribute of an existing content item. \sked{K14} \\
\addlinespace
REQ-F-009 & The system shall allow deleting a content item from the catalog. A deleted item is automatically removed from all collections. \sked{K15} \\
\addlinespace
REQ-F-024 & The add and edit item form shall display type-specific optional fields depending on the selected type (see \tabref{tab:type-fields}). \sked{K30} \\
\addlinespace
REQ-F-025 & Each content item shall have a cover image: by default --- one of three illustrations corresponding to the type. The user may replace it with their own image (optional). \sked{K31} \\
\addlinespace
REQ-F-028 & Cover image formats: JPEG, PNG, WebP; maximum file size: 5~MB. The system notifies the user when limits are violated. \sked{K48} \\
\bottomrule
\end{tabular}
\end{table}

\begin{table}[h]
\caption{Type-Specific Content Item Fields}
\label{tab:type-fields}
\begin{tabular}{L{3cm} L{11cm}}
\toprule
\textbf{Type} & \textbf{Optional Fields} \\
\midrule
Book & Author, number of pages \\
TV Series & Number of seasons \\
Film & Director, duration \\
\bottomrule
\end{tabular}
\end{table}
\sked{K30}

\subsection{Tags}

\begin{table}[h]
\caption{Tag Requirements}
\label{tab:tehy}
\begin{tabular}{L{2.5cm} L{11.5cm}}
\toprule
\textbf{ID} & \textbf{Requirement} \\
\midrule
REQ-F-010 & The system shall allow adding arbitrary tags to a content item. When entering a tag, the system offers autocomplete based on the user's previously entered tags. One item can have multiple tags. A tag can be removed at any time. \sked{K16} \\
\addlinespace
REQ-F-011 & The system shall provide catalog filtering and search by tag. \sked{K17} \\
\bottomrule
\end{tabular}
\end{table}

\subsection{Sessions (Repeated Viewings / Readings)}

\begin{table}[h]
\caption{Session Requirements}
\label{tab:sessions}
\begin{tabular}{L{2.5cm} L{11.5cm}}
\toprule
\textbf{ID} & \textbf{Requirement} \\
\midrule
REQ-F-012 & The system shall allow recording multiple Sessions for a single content item; each Session has its own date, rating, and note. \sked{K18} \\
\addlinespace
REQ-F-013 & The system shall display the list of all Sessions of a content item in reverse chronological order (newest to oldest), with the total count. \sked{K19} \\
\addlinespace
REQ-F-014 & The system shall allow editing and deleting an individual Session. \sked{K20} \\
\bottomrule
\end{tabular}
\end{table}

\subsection{Collections}

\begin{table}[h]
\caption{Collection Requirements}
\label{tab:dobirkы}
\begin{tabular}{L{2.5cm} L{11.5cm}}
\toprule
\textbf{ID} & \textbf{Requirement} \\
\midrule
REQ-F-015 & The system shall allow adding a single catalog item to multiple collections simultaneously (many-to-many relationship). \sked{K21} \\
\addlinespace
REQ-F-016 & The system shall allow creating a new named collection. \sked{K22} \\
\addlinespace
REQ-F-017 & The system shall allow editing a collection name. \sked{K23} \\
\addlinespace
REQ-F-018 & The system shall allow deleting a collection. Deleting a collection does not delete items from the catalog. \sked{K24} \\
\addlinespace
REQ-F-019 & The system shall allow removing an item from a collection without deleting it from the catalog or other collections. \sked{K25} \\
\addlinespace
REQ-F-020 & The system shall display the contents of a collection --- the list of items it contains, with the name, type, and status of each. \sked{K26} \\
\addlinespace
REQ-F-027 & The system shall provide sorting of items within a collection: by date added (newest~/ oldest) and by rating (highest to lowest and vice versa). \sked{K33} \\
\bottomrule
\end{tabular}
\end{table}

\subsection{Automatic Completion Date Fill}

\begin{table}[h]
\caption{Completion Date Requirement}
\label{tab:data}
\begin{tabular}{L{2.5cm} L{11.5cm}}
\toprule
\textbf{ID} & \textbf{Requirement} \\
\midrule
REQ-F-021 & When changing an item's status to ``completed'', the completion date field shall be automatically filled with the current date if no date was entered manually. \sked{K27} \\
\bottomrule
\end{tabular}
\end{table}

\subsection{Personal Statistics}

\begin{table}[h]
\caption{Personal Statistics Requirements}
\label{tab:stat}
\begin{tabular}{L{2.5cm} L{11.5cm}}
\toprule
\textbf{ID} & \textbf{Requirement} \\
\midrule
REQ-F-022 & The system shall display personal statistics: the number of items with ``completed'' status broken down by type (books~/ TV series~/ films) and by selected time period. \sked{K28} \\
\addlinespace
REQ-F-022a & Time period selection for statistics --- from predefined options: ``this month'', ``this year'', ``all time''. \sked{K51} \\
\bottomrule
\end{tabular}
\end{table}

\section{Non-Functional Requirements}
\label{sec:nefunktsionalni-vymohy}

\subsection{Performance}

\textbf{REQ-NF-001.} System pages shall load in an acceptable time for the user under normal network connection conditions. \sked{K34}

\subsection{Security}

\textbf{REQ-NF-002.} User passwords are stored exclusively in a protected form --- as hashes (e.g., bcrypt). Access to protected resources requires an active authenticated session. \sked{K35}

\textbf{REQ-NF-007.} The authorization token is stored in an HTTP-only cookie to prevent XSS attacks. All requests that modify system state must be protected against CSRF attacks. \sked{K47}

\subsection{Reliability}

\textbf{REQ-NF-003.} All user data is stored in the database. No user action shall result in the loss of previously entered data without explicit deletion confirmation from the user. \sked{K36}

\subsection{Usability}

\textbf{REQ-NF-004.} The system's interface shall be responsive for desktop and mobile browsers. Separate tablet optimization is not planned for version~v1. \sked{K37}

\textbf{REQ-NF-008.} When the catalog is empty, a collection is empty, or search~/ filter results are absent, the system shall display an informative message instead of an empty list. \sked{K52}

\subsection{Compatibility}

\textbf{REQ-NF-005.} The system shall function correctly in the latest two versions of: Google~Chrome, Mozilla~Firefox, Microsoft~Edge, and Apple~Safari (including mobile Safari on iOS). \sked{K45}

\subsection{Input Validation}

\textbf{REQ-NF-006.} During new user registration, the system shall validate:
\begin{itemize}
  \item email --- format per RFC~5322;
  \item password --- minimum length of 8 characters;
  \item nickname --- 2 to 50 characters.
\end{itemize}
If any constraint is violated, the system displays a specific error message indicating the reason for rejection. \sked{K46}

\subsection{Development Process}

\textbf{REQ-NF-009.} After each test run, the development agent shall write a log file to the \texttt{logs/} directory. The log file must contain: total number of tests, passed~/ failed~/ skipped count, status of each individual test (name~+ result), and execution time of the test suite. \sked{K54}

\section{Acceptance Criteria}
\label{sec:acceptance-criteria}

\begin{longtable}{L{2cm} L{2.5cm} L{9.5cm}}
\caption{Requirement Acceptance Criteria} \label{tab:ac} \\
\toprule
\textbf{AC ID} & \textbf{Req} & \textbf{Acceptance Criterion} \\
\midrule
\endfirsthead
\toprule
\textbf{AC ID} & \textbf{Req} & \textbf{Acceptance Criterion} \\
\midrule
\endhead
\midrule
\multicolumn{3}{r}{\emph{Continued on next page}} \\
\endfoot
\bottomrule
\endlastfoot

AC-001 & REQ-F-001 & After filling in a valid email, password, and nickname and confirming registration, the user gains access to an empty personal catalog. Re-registration with the same email is not possible: the system displays a message about an existing account. \sked{K7} \\
\addlinespace
AC-002 & REQ-F-002 & After entering the correct email and password, the user is redirected to their catalog. With incorrect credentials --- an error message is displayed and login does not proceed. \sked{K8} \\
\addlinespace
AC-003 & REQ-F-003 & After logout, accessing protected pages redirects to the login page. Re-authentication restores access. \sked{K9} \\
\addlinespace
AC-004 & REQ-F-004 & The catalog displays all added items; for each item, the title, type, status, and rating are visible. \sked{K10} \\
\addlinespace
AC-005 & REQ-F-005 & When selecting type ``Books'', only items of type ``book'' are displayed; when the filter is cleared --- all items are shown again. \sked{K11} \\
\addlinespace
AC-006 & REQ-F-006 & When selecting status ``completed'', only items with that status are displayed; when cleared --- all items are shown again. \sked{K12} \\
\addlinespace
AC-007 & REQ-F-007 & After filling in required fields (title, type), the item appears in the catalog with its corresponding attributes and cover image. \sked{K13} \\
\addlinespace
AC-008 & REQ-F-008 & After editing and saving, the updated attributes are shown in the catalog and on the item page. \sked{K14} \\
\addlinespace
AC-009 & REQ-F-009 & After confirming deletion, the item disappears from the catalog and from all collections it belonged to. \sked{K15} \\
\addlinespace
AC-010 & REQ-F-010 & An added tag appears on the item card. When typing in the tag field --- autocomplete suggestions from previously used tags are shown. Removing a tag does not affect other attributes. \sked{K16} \\
\addlinespace
AC-011 & REQ-F-011 & When filtering by tag, only items with that tag are displayed; the result updates when the filter is changed or cleared. \sked{K17} \\
\addlinespace
AC-012 & REQ-F-012 & For an item with two or more Sessions --- a list of all records with date, rating, and note is shown. \sked{K18} \\
\addlinespace
AC-013 & REQ-F-013 & Sessions are sorted from newest to oldest; the total count is indicated. \sked{K19} \\
\addlinespace
AC-014 & REQ-F-014 & After editing a Session --- updated data is shown. After deletion --- the session disappears and the count decreases. \sked{K20} \\
\addlinespace
AC-015 & REQ-F-015 & An item is simultaneously displayed in each collection it has been added to. Removal from one collection does not affect the others. \sked{K21} \\
\addlinespace
AC-016 & REQ-F-016 & After entering a name, the new empty collection appears in the collections list. \sked{K22} \\
\addlinespace
AC-017 & REQ-F-017 & After saving, the new name is displayed in the list and in the page heading. \sked{K23} \\
\addlinespace
AC-018 & REQ-F-018 & After confirming deletion, the collection disappears; items remain in the catalog. \sked{K24} \\
\addlinespace
AC-019 & REQ-F-019 & After removing an item from a collection, it remains in the catalog and in other collections. \sked{K25} \\
\addlinespace
AC-020 & REQ-F-020 & The collection page displays all items with title, type, and status. \sked{K26} \\
\addlinespace
AC-021 & REQ-F-021 & If no date is entered when changing status to ``completed'' --- the current date is automatically set. If entered manually --- the entered value is saved. \sked{K27} \\
\addlinespace
AC-022 & REQ-F-022 & The statistics page displays the count of completed items by type and by the selected time period. \sked{K28} \\
\addlinespace
AC-023 & REQ-F-023 & When entering part of a title, the catalog displays only matching items; when cleared --- all items are shown. \sked{K29} \\
\addlinespace
AC-024 & REQ-F-024 & The form displays type-specific fields only after a type is selected; all fields are optional. \sked{K30} \\
\addlinespace
AC-025 & REQ-F-025 & Without a custom image --- the default illustration for the type is shown. After upload --- the custom image is shown; it can be removed to restore the default. \sked{K31} \\
\addlinespace
AC-026 & REQ-F-026 & After selecting a sort parameter, the catalog is reordered; changing the parameter reorders it again. \sked{K32} \\
\addlinespace
AC-027 & REQ-F-027 & After selecting a sort parameter within a collection, the list is reordered. \sked{K33} \\
\addlinespace
AC-028 & REQ-F-028 & An attempt to upload a file of an unsupported format or larger than 5~MB --- an error message is shown and the file is not saved. \sked{K48} \\
\addlinespace
AC-029 & REQ-F-029 & After changing the nickname in account settings, the new name is displayed in the interface. \sked{K49} \\
\addlinespace
AC-030 & REQ-F-022a & The statistics page contains a time period selector: ``this month'', ``this year'', ``all time''. \sked{K51} \\
\addlinespace
AC-031 & REQ-NF-006 & When a password shorter than 8 characters is entered --- a specific message is shown. With an invalid email or a nickname outside 2--50 characters --- a corresponding message is shown. \sked{K46} \\
\addlinespace
AC-032 & REQ-NF-008 & When the catalog is empty, a collection is empty, or search results are absent, an informative message is displayed (not an empty list). \sked{K52} \\
\addlinespace
AC-033 & REQ-NF-009 & After running tests, a file with the results of the last run exists in the \texttt{logs/} directory. The file contains: total number of tests, count of passed~/ failed~/ skipped, status of each test (name~+ result), and execution time of the suite. \sked{K54} \\

\end{longtable}

\section{BDD Scenarios}
\label{sec:bdd}

Below are BDD scenarios in Gherkin format for key system functions.

\begin{codelisting}
\caption{New User Registration}
\label{lst:bdd}
\begin{skedcode}
Scenario: New user registration
  Given the registration page is open
  And email "test@example.com" is not yet registered in the system
  When the user enters email "test@example.com", a password, and a nickname
    and confirms registration
  Then the user gains access to an empty personal catalog
  And a repeated registration attempt with the same email shows the message
    "An account with this email already exists"
\end{skedcode}
\end{codelisting}
\sked{K7}

\begin{codelisting}
\caption{Account Login}
\label{lst:bdd-login}
\begin{skedcode}
Scenario: Login with correct credentials
  Given the user is registered in the system
  When they enter the correct email and password and click "Sign in"
  Then they are redirected to their personal catalog
  And they see their previously added items
\end{skedcode}
\end{codelisting}
\sked{K8}

\begin{codelisting}
\caption{Repeated Viewing with a Different Rating}
\label{lst:bdd-session}
\begin{skedcode}
Scenario: Re-watching a film with a different rating
  Given content item "Film X" already has one Session with rating 9
  When the user adds a new Session for the same item
    with rating 5 and note "didn't enjoy it the second time"
  Then the content item has two Sessions, each with its own date,
    rating, and note
  And both ratings are stored separately, neither overwrites the other
\end{skedcode}
\end{codelisting}
\sked{K18}

\begin{codelisting}
\caption{Item in Multiple Collections}
\label{lst:bdd-collection}
\begin{skedcode}
Scenario: An item belongs to multiple collections simultaneously
  Given the user has item "Book Y" in the catalog
  And collections "Favourites" and "2026" exist
  When the user adds "Book Y" to both collections
  Then "Book Y" is displayed in the contents of both collections
  And removing "Book Y" from collection "2026" does not remove it
    from "Favourites" or from the catalog
\end{skedcode}
\end{codelisting}
\sked{K21,K25}

\begin{codelisting}
\caption{Automatic Completion Date Fill}
\label{lst:bdd-date}
\begin{skedcode}
Scenario: Automatic completion date fill
  Given the user is editing an item with status "in progress"
  When the user changes the status to "completed"
    and does not enter a date manually
  Then the system automatically sets the completion date
    equal to the current date
\end{skedcode}
\end{codelisting}
\sked{K27}

\begin{codelisting}
\caption{Filtering by Tag}
\label{lst:bdd-tag}
\begin{skedcode}
Scenario: Filtering the catalog by tag
  Given the catalog contains three items: two tagged "fantasy" and one without
  When the user selects the filter by tag "fantasy"
  Then only the two items tagged with that tag are displayed
  And when the filter is cleared, all three items are displayed again
\end{skedcode}
\end{codelisting}
\sked{K17}

\section{Requirements Traceability Matrix}
\label{sec:traceability}

\begin{longtable}{L{2.8cm} L{3.2cm} L{2.5cm} L{5.5cm}}
\caption{Traceability Matrix} \label{tab:traceability} \\
\toprule
\textbf{REQ ID} & \textbf{AC ID} & \textbf{BDD} & \textbf{Note} \\
\midrule
\endfirsthead
\toprule
\textbf{REQ ID} & \textbf{AC ID} & \textbf{BDD} & \textbf{Note} \\
\midrule
\endhead
\midrule
\multicolumn{4}{r}{\emph{Continued on next page}} \\
\endfoot
\bottomrule
\endlastfoot

REQ-F-001 & AC-001 & \lstref{lst:bdd} & Registration \\
REQ-F-002 & AC-002 & \lstref{lst:bdd-login} & Login \\
REQ-F-003 & AC-003 & --- & Logout \\
REQ-F-004 & AC-004 & --- & Catalog \\
REQ-F-005 & AC-005 & --- & Filter by type \\
REQ-F-006 & AC-006 & --- & Filter by status \\
REQ-F-007 & AC-007 & --- & Add item \\
REQ-F-008 & AC-008 & --- & Edit item \\
REQ-F-009 & AC-009 & --- & Delete item \\
REQ-F-010 & AC-010 & --- & Tags with autocomplete \\
REQ-F-011 & AC-011 & \lstref{lst:bdd-tag} & Filter by tag \\
REQ-F-012 & AC-012 & \lstref{lst:bdd-session} & Sessions \\
REQ-F-013 & AC-013 & --- & View Sessions \\
REQ-F-014 & AC-014 & --- & Edit~/ delete Session \\
REQ-F-015 & AC-015 & \lstref{lst:bdd-collection} & Many-to-many collections \\
REQ-F-016 & AC-016 & --- & Create collection \\
REQ-F-017 & AC-017 & --- & Edit collection \\
REQ-F-018 & AC-018 & --- & Delete collection \\
REQ-F-019 & AC-019 & \lstref{lst:bdd-collection} & Remove from collection \\
REQ-F-020 & AC-020 & --- & View collection \\
REQ-F-021 & AC-021 & \lstref{lst:bdd-date} & Auto completion date \\
REQ-F-022 & AC-022 & --- & Statistics \\
REQ-F-022a & AC-030 & --- & Statistics time period \\
REQ-F-023 & AC-023 & --- & Search by title \\
REQ-F-024 & AC-024 & --- & Type-specific fields \\
REQ-F-025 & AC-025 & --- & Cover image \\
REQ-F-026 & AC-026 & --- & Catalog sorting \\
REQ-F-027 & AC-027 & --- & Collection sorting \\
REQ-F-028 & AC-028 & --- & Cover image constraints \\
REQ-F-029 & AC-029 & --- & Profile editing \\
REQ-NF-001 & --- & --- & Performance \\
REQ-NF-002 & --- & --- & Security~/ hashing \\
REQ-NF-003 & --- & --- & Reliability \\
REQ-NF-004 & --- & --- & Responsive interface \\
REQ-NF-005 & --- & --- & Browser compatibility \\
REQ-NF-006 & AC-031 & --- & Input validation \\
REQ-NF-007 & --- & --- & Session security (CSRF, HTTP-only) \\
REQ-NF-008 & AC-032 & --- & Empty states \\
REQ-NF-009 & AC-033 & --- & Test log files \\

\end{longtable}


\skedappendix

\section{Features Outside the MVP Scope}
\label{sec:out-of-scope}

The following features are deliberately excluded from version~v1 to constrain the scope of the first release. They may be considered for future versions of the system. \sked{K43,K50}

\begin{itemize}
  \item Offline mode --- full operation without an internet connection.
  \item Push and email notifications.
  \item Multi-language interface.
  \item ``Friends'' feature --- mutual subscriptions, shared access to collections.
  \item Recommendations block or ``Popular among users''.
  \item Integration with external media content databases (TMDB, Google Books, etc.).
  \item Administrative panel.
  \item Email confirmation at registration (sending a verification code by email).
  \item Changing email address and changing password in account settings.
  \item Account deletion by the user.
\end{itemize}


\skedbibliography

\end{document}
